% ============================================================
%  SAMURAI — Теория автоматического управления
%  ОДУ-модель, модальный/оптимальный регулятор, MPC, наблюдатель
%  Главный файл документа
%  Компиляция: pdflatex main.tex   (два прохода для оглавления)
% ============================================================
\documentclass[12pt, a4paper, openany]{book}

% ── Кодировка и язык ─────────────────────────────────────────────
\usepackage[utf8]{inputenc}
\usepackage[T2A]{fontenc}
\usepackage[russian, english]{babel}

% ── Геометрия страницы ───────────────────────────────────────────
\usepackage[
  top=25mm, bottom=25mm,
  left=30mm, right=20mm
]{geometry}

% ── Математика ───────────────────────────────────────────────────
\usepackage{amsmath, amssymb, amsthm}
\usepackage{mathtools}
\usepackage{bm}

% ── Графика и цвет ───────────────────────────────────────────────
\usepackage{xcolor}
\usepackage{graphicx}
\usepackage{tikz}
\usetikzlibrary{arrows.meta, positioning, shapes.geometric, calc}

% ── Код ──────────────────────────────────────────────────────────
\usepackage{listings}
\usepackage{courier}

\definecolor{codebg}{RGB}{245,247,250}
\definecolor{codeframe}{RGB}{180,190,210}
\definecolor{keyword}{RGB}{0,100,200}
\definecolor{comment}{RGB}{100,130,100}
\definecolor{string}{RGB}{163,21,21}
\definecolor{number}{RGB}{9,134,88}

% UTF-8 support for listings: literate replacement for Cyrillic chars.
% Без этого pdflatex падает на байтах UTF-8 в комментариях кода.
\lstset{
  extendedchars=true,
  inputencoding=utf8,
  literate=
    {а}{{\selectfont\char224}}1{б}{{\selectfont\char225}}1
    {в}{{\selectfont\char226}}1{г}{{\selectfont\char227}}1
    {д}{{\selectfont\char228}}1{е}{{\selectfont\char229}}1
    {ё}{{\selectfont\char184}}1{ж}{{\selectfont\char230}}1
    {з}{{\selectfont\char231}}1{и}{{\selectfont\char232}}1
    {й}{{\selectfont\char233}}1{к}{{\selectfont\char234}}1
    {л}{{\selectfont\char235}}1{м}{{\selectfont\char236}}1
    {н}{{\selectfont\char237}}1{о}{{\selectfont\char238}}1
    {п}{{\selectfont\char239}}1{р}{{\selectfont\char240}}1
    {с}{{\selectfont\char241}}1{т}{{\selectfont\char242}}1
    {у}{{\selectfont\char243}}1{ф}{{\selectfont\char244}}1
    {х}{{\selectfont\char245}}1{ц}{{\selectfont\char246}}1
    {ч}{{\selectfont\char247}}1{ш}{{\selectfont\char248}}1
    {щ}{{\selectfont\char249}}1{ъ}{{\selectfont\char250}}1
    {ы}{{\selectfont\char251}}1{ь}{{\selectfont\char252}}1
    {э}{{\selectfont\char253}}1{ю}{{\selectfont\char254}}1
    {я}{{\selectfont\char255}}1
    {А}{{\selectfont\char192}}1{Б}{{\selectfont\char193}}1
    {В}{{\selectfont\char194}}1{Г}{{\selectfont\char195}}1
    {Д}{{\selectfont\char196}}1{Е}{{\selectfont\char197}}1
    {Ё}{{\selectfont\char168}}1{Ж}{{\selectfont\char198}}1
    {З}{{\selectfont\char199}}1{И}{{\selectfont\char200}}1
    {Й}{{\selectfont\char201}}1{К}{{\selectfont\char202}}1
    {Л}{{\selectfont\char203}}1{М}{{\selectfont\char204}}1
    {Н}{{\selectfont\char205}}1{О}{{\selectfont\char206}}1
    {П}{{\selectfont\char207}}1{Р}{{\selectfont\char208}}1
    {С}{{\selectfont\char209}}1{Т}{{\selectfont\char210}}1
    {У}{{\selectfont\char211}}1{Ф}{{\selectfont\char212}}1
    {Х}{{\selectfont\char213}}1{Ц}{{\selectfont\char214}}1
    {Ч}{{\selectfont\char215}}1{Ш}{{\selectfont\char216}}1
    {Щ}{{\selectfont\char217}}1{Ъ}{{\selectfont\char218}}1
    {Ы}{{\selectfont\char219}}1{Ь}{{\selectfont\char220}}1
    {Э}{{\selectfont\char221}}1{Ю}{{\selectfont\char222}}1
    {Я}{{\selectfont\char223}}1
    {—}{{\textemdash}}1{–}{{\textendash}}1{«}{{\guillemotleft}}1
    {»}{{\guillemotright}}1{№}{{N\textsuperscript{o}}}1
    {²}{{\textsuperscript{2}}}1{·}{{$\cdot$}}1
}

\lstdefinestyle{pythonstyle}{
  language=Python,
  backgroundcolor=\color{codebg},
  frame=single,
  rulecolor=\color{codeframe},
  framesep=4pt,
  basicstyle=\footnotesize\ttfamily,
  keywordstyle=\color{keyword}\bfseries,
  commentstyle=\color{comment}\itshape,
  stringstyle=\color{string},
  numberstyle=\scriptsize\ttfamily\color{gray},
  numbers=left,
  numbersep=10pt,
  xleftmargin=2.5em,
  showstringspaces=false,
  breaklines=true,
  tabsize=4,
  captionpos=b,
  morekeywords={True, False, None, self},
}

\lstdefinestyle{matlabstyle}{
  language=Matlab,
  backgroundcolor=\color{codebg},
  frame=single,
  rulecolor=\color{codeframe},
  framesep=4pt,
  basicstyle=\footnotesize\ttfamily,
  keywordstyle=\color{keyword}\bfseries,
  commentstyle=\color{comment}\itshape,
  stringstyle=\color{string},
  numberstyle=\scriptsize\ttfamily\color{gray},
  numbers=left,
  numbersep=10pt,
  xleftmargin=2.5em,
  showstringspaces=false,
  breaklines=true,
  tabsize=2,
  captionpos=b,
}

\lstset{style=pythonstyle}

% ── Таблицы ──────────────────────────────────────────────────────
\usepackage{booktabs}
\usepackage{array}
\usepackage{longtable}
\usepackage{float}

% ── Гиперссылки ──────────────────────────────────────────────────
\usepackage[
  colorlinks=true,
  linkcolor=blue!70!black,
  citecolor=green!60!black,
  urlcolor=blue!80!black,
  bookmarks=true,
  bookmarksnumbered=true
]{hyperref}

% ── Колонтитулы ──────────────────────────────────────────────────
\usepackage{fancyhdr}
\pagestyle{fancy}
\fancyhf{}
\fancyhead[LE,RO]{\small\thepage}
\fancyhead[LO]{\small\nouppercase{\rightmark}}
\fancyhead[RE]{\small\nouppercase{\leftmark}}
\renewcommand{\headrulewidth}{0.4pt}

% ── Параграфы ────────────────────────────────────────────────────
\usepackage{parskip}
\setlength{\parindent}{0pt}

% ── Окружения ────────────────────────────────────────────────────
\usepackage{tcolorbox}
\tcbuselibrary{skins, breakable}

\newtcolorbox{formulabox}[1][]{
  enhanced, breakable,
  colback=blue!3!white, colframe=blue!50!black,
  fonttitle=\bfseries, title=Формула, #1
}
\newtcolorbox{notebox}[1][]{
  enhanced,
  colback=orange!8!white, colframe=orange!70!black,
  fonttitle=\bfseries, title=Примечание, #1
}
\newtcolorbox{filebox}[1][]{
  enhanced,
  colback=gray!6!white, colframe=gray!50!black,
  fonttitle=\bfseries\ttfamily, #1
}
\newtcolorbox{warnbox}[1][]{
  enhanced,
  colback=red!4!white, colframe=red!70!black,
  fonttitle=\bfseries, title=Внимание, #1
}

% ── Нумерация уравнений по главам ────────────────────────────────
\numberwithin{equation}{chapter}

% ── Макросы ──────────────────────────────────────────────────────
\newcommand{\filepath}[1]{\texttt{\small\color{gray!80!black}#1}}
\newcommand{\vect}[1]{\bm{#1}}
\newcommand{\mat}[1]{\mathbf{#1}}
\newcommand{\T}{^{\mathsf{T}}}
\DeclareMathOperator*{\argmin}{arg\,min}

% ============================================================
\begin{document}

% ── Титульная страница ───────────────────────────────────────────
\begin{titlepage}
\centering
\vspace*{2cm}
{\Huge\bfseries SAMURAI Robot}\\[0.4cm]
{\Large\bfseries Теория автоматического управления}\\[0.3cm]
{\large ОДУ-модель, синтез регуляторов, MPC, наблюдатель состояния}\\[1cm]
\rule{\textwidth}{1pt}\\[0.5cm]
{\large Технический справочник по регулятору движения робота}\\[0.5cm]
\rule{\textwidth}{1pt}\\[2cm]

\begin{tabular}{ll}
\textbf{Объект управления:} & 2-моторный гусеничный робот (Adeept HAT)\\
\textbf{Размерность состояния:} & $n=5$ (координаты + скорости)\\
\textbf{Размерность управления:} & $r=2$ (linear, angular)\\
\textbf{Период дискретизации:} & $T_s = 0.05$ с (20 Гц)\\
\textbf{Метод линеаризации:} & ряд Тейлора в опорной точке\\
\textbf{Синтез регуляторов:} & модальный, ЛКР, MPC\\
\textbf{Наблюдатель:} & полноразмерный (Luenberger)\\
\textbf{Реализация:} & MATLAB (offline) $\to$ YAML $\to$ Python (онлайн)\\
\end{tabular}

\vfill
{\large Источник теории: Козлов В.\,Н., Куприянов В.\,Е., Шашихин В.\,Н.\\
``Теория автоматического управления'', СПбГПУ, 2008}\\[0.5cm]
{\large\today}
\end{titlepage}

% ── Оглавление ───────────────────────────────────────────────────
\tableofcontents
\newpage

% ── Введение ─────────────────────────────────────────────────────
\chapter*{Введение}
\addcontentsline{toc}{chapter}{Введение}

Документ содержит полное математическое описание системы автоматического
управления (САУ) гусеничным роботом SAMURAI. Изложение построено по схеме
\textbf{«модели $\to$ анализ $\to$ синтез»} согласно учебному пособию
\textit{Козлов В.\,Н., Куприянов В.\,Е., Шашихин В.\,Н. ``Теория автоматического
управления''} (СПбГПУ, 2008).

\section*{Постановка задачи}

Робот SAMURAI имеет 2 задних мотора (M1, M2) с дифференциальной кинематикой
и должен отрабатывать заданные траектории движения. Существующая система
управления использует прямую трансляцию команд скорости (\texttt{cmd\_vel})
в ШИМ моторов через статические калибровочные коэффициенты, что:

\begin{itemize}
  \item не учитывает инерционность моторов ($\tau \approx 0.15$ с);
  \item не оптимизирует управление по интегральному критерию;
  \item не предсказывает поведение системы на горизонте.
\end{itemize}

Цель работы --- построить \textbf{линейную модель} робота, синтезировать
\textbf{ЛКР} и \textbf{MPC-регулятор}, и встроить их в код управления так,
чтобы параметры регулятора были \textbf{редактируемы из конфига}. Это
позволяет:

\begin{enumerate}
  \item улучшать поведение робота, подбирая матрицы $\mat{Q}$, $\mat{R}$;
  \item \textbf{намеренно ухудшать} поведение для отладки и демонстрации
        неустойчивости (выбор «неправильных» полюсов и весов);
  \item менять режим управления на лету (off/LQR/MPC) без перекомпиляции.
\end{enumerate}

\section*{Структура документа}

\begin{center}
\begin{tabular}{ll}
\toprule
\textbf{Глава} & \textbf{Содержание}\\
\midrule
1. Модель робота & Нелинейная ОДУ, линеаризация, дискретизация\\
2. Синтез регуляторов & Модальный, ЛКР через Риккати\\
3. MPC & Прогнозирующее управление, QP-задача\\
4. Реализация в коде & MATLAB $\to$ YAML $\to$ Python\\
\bottomrule
\end{tabular}
\end{center}

\section*{Используемые обозначения}

\begin{center}
\begin{tabular}{ll}
\toprule
$\vect{x}(t) \in \mathbb{R}^n$ & вектор состояния объекта\\
$\vect{u}(t) \in \mathbb{R}^r$ & вектор управляющих воздействий\\
$\vect{y}(t) \in \mathbb{R}^m$ & вектор наблюдаемых выходов\\
$\mat{A}, \mat{B}, \mat{C}, \mat{D}$ & матрицы модели в пространстве состояний\\
$\mat{K}$ & матрица обратной связи (регулятор)\\
$\mat{L}$ & матрица обратной связи наблюдателя\\
$\mat{Q}, \mat{R}$ & весовые матрицы критерия качества\\
$\mat{P}$ & решение уравнения Риккати\\
$T_s$ & период дискретизации\\
$N$ & горизонт прогноза MPC\\
\bottomrule
\end{tabular}
\end{center}

% ── Главы ────────────────────────────────────────────────────────
% ============================================================
%  Глава 1 — Математическая модель робота SAMURAI
% ============================================================
\chapter{Математическая модель робота SAMURAI}
\label{ch:model}

\section{Физическая постановка}
\label{sec:physical}

Робот SAMURAI --- гусеничное шасси с дифференциальным приводом.
Контроллер PCA9685 управляет двумя задними моторами (M1 левый, M2 правый);
передние моторы физически отсутствуют (см. \filepath{hardware\_presets/samurai\_default.yaml}).

\begin{filebox}[title={Параметры объекта (\filepath{config.yaml})}]
\begin{verbatim}
simulator:
  wheel_base:    0.17    # m  (track separation)
  max_linear:    0.30    # m/s
  max_angular:   2.00    # rad/s
\end{verbatim}
\end{filebox}

\textbf{Команда движения} формируется в виде вектора $\vect{u} = (u_v, u_\omega)\T$,
где $u_v$ --- желаемая линейная скорость [м/с], $u_\omega$ --- желаемая
угловая скорость [рад/с]. Дифференциальная кинематика преобразует $\vect{u}$
в скорости гусениц:
\begin{equation}
  v_L = u_v + \tfrac{B}{2}\,u_\omega, \qquad
  v_R = u_v - \tfrac{B}{2}\,u_\omega,
  \label{eq:diff_kinematics}
\end{equation}
где $B$ --- ширина базы (\texttt{wheel\_base}).

\section{Нелинейная ОДУ модели}
\label{sec:ode}

В отличие от тривиальной кинематической модели, реальный робот имеет
\textbf{инерционность моторов}: команда $u_v$ не отрабатывается мгновенно,
а через апериодическое звено первого порядка с постоянной времени $\tau_v$.
Аналогично для угловой скорости (постоянная $\tau_\omega$).

Введём \textbf{вектор состояния} ($n=5$):
\begin{equation}
  \vect{x}(t) = \big(\,p_x,\; p_y,\; \theta,\; v,\; \omega\,\big)\T
  \label{eq:state_vector}
\end{equation}

\begin{itemize}
  \item $p_x, p_y$ --- координаты центра робота в мировой системе [м];
  \item $\theta$ --- угол курса (yaw) [рад];
  \item $v$ --- фактическая линейная скорость [м/с];
  \item $\omega$ --- фактическая угловая скорость [рад/с].
\end{itemize}

\begin{formulabox}[title={Нелинейная модель в форме $\dot{\vect{x}} = \vect{f}(\vect{x}, \vect{u})$}]
\begin{equation}
  \begin{cases}
    \dot{p}_x      = v\cos\theta,\\[2pt]
    \dot{p}_y      = v\sin\theta,\\[2pt]
    \dot{\theta}   = \omega,\\[2pt]
    \dot{v}        = \dfrac{1}{\tau_v}\,(u_v - v),\\[6pt]
    \dot{\omega}   = \dfrac{1}{\tau_\omega}\,(u_\omega - \omega).
  \end{cases}
  \label{eq:nonlinear_ode}
\end{equation}
\end{formulabox}

Первые три уравнения --- кинематика дифференциального робота
(идентична уравнениям 1.1 из учебника Козлова, формулы для $\dot{x}$ и $\dot{z}$).
Последние два --- динамика моторов как RC-звено с постоянной времени $\tau$,
определяемой механической инерцией ротора и противо-ЭДС.

\begin{notebox}
Постоянные времени $\tau_v \approx 0{,}15$ с, $\tau_\omega \approx 0{,}10$ с
оценены экспериментально по графику разгона (см.
\filepath{tests/test\_motor\_lag.py}). Для других шасси значения нужно
переопределить --- они вынесены в \texttt{config.yaml} раздел \texttt{control.plant}.
\end{notebox}

\section{Линеаризация в опорной точке}
\label{sec:linearization}

Выберем \textbf{опорный режим} --- равномерное прямолинейное движение
вперёд со скоростью $v_0$:
\begin{equation}
  \bar{\vect{x}} = (\bar{p}_x,\; \bar{p}_y,\; 0,\; v_0,\; 0)\T,
  \qquad
  \bar{\vect{u}} = (v_0,\; 0)\T.
  \label{eq:operating_point}
\end{equation}

Подстановка $\bar{\vect{x}}, \bar{\vect{u}}$ в \eqref{eq:nonlinear_ode}
даёт правые части $\dot{\bar{\vect{x}}} = (v_0,\,0,\,0,\,0,\,0)\T$, что
соответствует нулевым отклонениям состояния от опорной траектории при
постоянном $v_0$.

Согласно методу линеаризации (см. учебник п.~1.1, формулы 1.2),
уравнение для отклонений $\Delta\vect{x} = \vect{x} - \bar{\vect{x}}$,
$\Delta\vect{u} = \vect{u} - \bar{\vect{u}}$ имеет вид:
\begin{equation}
  \Delta\dot{\vect{x}}(t) = \mat{A}\,\Delta\vect{x}(t) + \mat{B}\,\Delta\vect{u}(t),
  \qquad
  \mat{A} = \left.\frac{\partial \vect{f}}{\partial\vect{x}}\right|_{\bar{\vect{x}},\bar{\vect{u}}},
  \quad
  \mat{B} = \left.\frac{\partial \vect{f}}{\partial\vect{u}}\right|_{\bar{\vect{x}},\bar{\vect{u}}}.
  \label{eq:linearization}
\end{equation}

\subsection{Вычисление матриц Якоби}

Распишем частные производные функций $f_i$ по $x_j$:
\begin{align*}
  \frac{\partial f_1}{\partial \theta} &= -v\sin\theta\big|_0 = 0, &
  \frac{\partial f_1}{\partial v} &= \cos\theta\big|_0 = 1,\\
  \frac{\partial f_2}{\partial \theta} &= v\cos\theta\big|_0 = v_0, &
  \frac{\partial f_2}{\partial v} &= \sin\theta\big|_0 = 0,\\
  \frac{\partial f_3}{\partial \omega} &= 1, &
  \frac{\partial f_4}{\partial v} &= -1/\tau_v,\\
  \frac{\partial f_5}{\partial \omega} &= -1/\tau_\omega.
\end{align*}

Остальные элементы матрицы $\mat{A}$ равны нулю. Получаем:

\begin{formulabox}[title={Линеаризованная модель в пространстве состояний}]
\begin{equation}
  \mat{A} = \begin{pmatrix}
    0 & 0 & 0    & 1 & 0 \\
    0 & 0 & v_0  & 0 & 0 \\
    0 & 0 & 0    & 0 & 1 \\
    0 & 0 & 0    & -1/\tau_v & 0 \\
    0 & 0 & 0    & 0 & -1/\tau_\omega
  \end{pmatrix},
  \qquad
  \mat{B} = \begin{pmatrix}
    0 & 0 \\
    0 & 0 \\
    0 & 0 \\
    1/\tau_v & 0 \\
    0 & 1/\tau_\omega
  \end{pmatrix}.
  \label{eq:AB_matrices}
\end{equation}
\end{formulabox}

Матрица наблюдения $\mat{C}$ выбирается по доступным датчикам.
В роботе SAMURAI прямо измеряются $p_x, p_y$ (одометрия),
$\theta$ (IMU), а $v, \omega$ восстанавливаются из последовательных
позиций. Будем считать наблюдаемыми все 5 координат:
\begin{equation}
  \mat{C} = \mat{I}_5,
  \qquad
  \mat{D} = \mat{0}_{5\times 2}.
  \label{eq:CD_matrices}
\end{equation}

\subsection{Численные значения}
\label{sec:numerics}

При $v_0 = 0{,}2$ м/с (профиль \texttt{normal}), $\tau_v = 0{,}15$ с,
$\tau_\omega = 0{,}10$ с:
\begin{equation}
  \mat{A} = \begin{pmatrix}
    0 & 0 & 0   & 1 & 0\\
    0 & 0 & 0.2 & 0 & 0\\
    0 & 0 & 0   & 0 & 1\\
    0 & 0 & 0   & -6.667 & 0\\
    0 & 0 & 0   & 0 & -10
  \end{pmatrix},
  \quad
  \mat{B} = \begin{pmatrix}
    0 & 0\\
    0 & 0\\
    0 & 0\\
    6.667 & 0\\
    0 & 10
  \end{pmatrix}.
  \label{eq:AB_numerical}
\end{equation}

\section{Анализ управляемости}
\label{sec:controllability}

Объект полностью управляем тогда и только тогда, когда матрица управляемости
имеет полный ранг (учебник, формула перед п.~3.1):
\begin{equation}
  \mat{S}_y = \big[\,\mat{B}\,\big|\,\mat{A}\mat{B}\,\big|\,\mat{A}^2\mat{B}\,\big|\,\ldots\,\big|\,\mat{A}^{n-1}\mat{B}\,\big],
  \qquad
  \operatorname{rank} \mat{S}_y = n.
  \label{eq:controllability}
\end{equation}

Для нашей системы $n=5$, $r=2$, поэтому $\mat{S}_y$ имеет размер $5 \times 10$.
Вычисление в MATLAB (\filepath{matlab/analyze\_system.m}, см. главу~\ref{ch:implementation})
показывает $\operatorname{rank}\mat{S}_y = 5$ при $v_0 \neq 0$ ---
объект полностью управляем.

\begin{warnbox}
При $v_0 = 0$ опорная точка вырождена: координата $p_y$ становится
неуправляемой (нельзя двигаться вбок без вращения). Для синтеза регулятора
необходимо $v_0 > 0$ или модель в окрестности нетривиальной траектории.
\end{warnbox}

\section{Дискретизация}
\label{sec:discretization}

Для реализации регулятора в реальном времени переходим к дискретной модели
с периодом $T_s = 0{,}05$ с (20 Гц --- частота цикла \texttt{motor\_node}).
Согласно п.~1.5 учебника, дискретизация ZOH (zero-order hold) даёт:
\begin{equation}
  \boxed{
    \mat{A}_d = e^{\mat{A} T_s}, \qquad
    \mat{B}_d = \int_0^{T_s} e^{\mat{A}\tau}\,d\tau \cdot \mat{B},
  }
  \label{eq:zoh_discretization}
\end{equation}
что эквивалентно вычислению блочной матричной экспоненты:
\begin{equation}
  \exp\!\left(
    \begin{pmatrix} \mat{A} & \mat{B} \\ \mat{0} & \mat{0} \end{pmatrix} T_s
  \right)
  =
  \begin{pmatrix} \mat{A}_d & \mat{B}_d \\ \mat{0} & \mat{I} \end{pmatrix}.
  \label{eq:block_exp}
\end{equation}

Дискретная модель:
\begin{equation}
  \vect{x}_{k+1} = \mat{A}_d\,\vect{x}_k + \mat{B}_d\,\vect{u}_k,
  \qquad
  \vect{y}_k = \mat{C}\,\vect{x}_k.
  \label{eq:discrete_model}
\end{equation}

\subsection{Соответствие собственных значений}

Собственные значения непрерывной и дискретной систем связаны соотношением
(п.~3.1 учебника):
\begin{equation}
  \tilde{\lambda}_i = e^{\lambda_i T_s}.
  \label{eq:eigval_mapping}
\end{equation}
Условия устойчивости:
\begin{itemize}
  \item непрерывная: $\operatorname{Re}\lambda_i < 0$ ---
        собственные значения в левой полуплоскости;
  \item дискретная: $|\tilde{\lambda}_i| < 1$ ---
        внутри единичного круга.
\end{itemize}

\begin{notebox}
Для непрерывного робота $\mat{A}$ имеет три нулевых собственных значения
(апериодические интеграторы по $p_x, p_y, \theta$) и два отрицательных:
$-1/\tau_v$, $-1/\tau_\omega$. Объект \textbf{не устойчив асимптотически}
(но и не неустойчив --- это интегратор), что и ожидаемо: робот без
управления продолжает двигаться по инерции.
\end{notebox}

% ============================================================
%  Глава 2 — Синтез регуляторов: модальный + ЛКР
% ============================================================
\chapter{Синтез регуляторов: модальный и оптимальный}
\label{ch:synthesis}

В этой главе строится статический регулятор $\vect{u} = -\mat{K}\,\vect{x}$,
размещающий полюса замкнутой системы в желаемых точках.
Излагаются два подхода (по разделам 3.1 и 3.3 учебника):
\textbf{модальный} (выбор полюсов вручную) и
\textbf{оптимальный ЛКР} (минимизация интегрального критерия).

\section{Постановка задачи стабилизации}
\label{sec:stabilization}

Дана линейная стационарная модель (по результатам главы~\ref{ch:model}):
\begin{equation}
  \dot{\vect{x}}(t) = \mat{A}\,\vect{x}(t) + \mat{B}\,\vect{u}(t),
  \qquad \vect{x}(0) = \vect{x}^0.
  \label{eq:plant_continuous}
\end{equation}

Требуется построить регулятор статической обратной связи:
\begin{equation}
  \boxed{\vect{u}(t) = -\mat{K}\,\vect{x}(t),}
  \label{eq:state_feedback}
\end{equation}
такой, что замкнутая система $\dot{\vect{x}} = (\mat{A}-\mat{B}\mat{K})\vect{x}$
асимптотически устойчива, и переходный процесс удовлетворяет требованиям
по длительности и перерегулированию.

\section{Модальный регулятор (pole placement)}
\label{sec:modal}

\subsection{Идея метода}

Зададим желаемые собственные значения замкнутой системы
$\{\lambda_i^{\text{уст}}\}_{i=1}^n$ и подберём $\mat{K}$ так,
чтобы спектр $\sigma(\mat{A}-\mat{B}\mat{K})$ совпал с этим множеством.

Согласно п.~3.1.3 учебника, для гарантии экспоненциального переходного
процесса заданной длительности $t_n$ выбираются вещественные собственные
значения, удовлетворяющие
\begin{equation}
  \min_i |\lambda_i^{\text{уст}}| \;\geq\; \frac{t_n}{3}.
  \label{eq:eigval_bound}
\end{equation}

\subsection{Каноническая форма Фробениуса}

Для одномерного входа ($r=1$) модальный регулятор синтезируется через
приведение к канонической форме (п.~3.1.2 учебника):
\begin{equation}
  \mat{A}_f = \begin{pmatrix}
    0 & 1 & 0 & \cdots & 0\\
    0 & 0 & 1 & \cdots & 0\\
    \vdots&&&\ddots&\vdots\\
    -a_n & -a_{n-1} & \cdots & & -a_1
  \end{pmatrix},
  \quad
  \vect{b}_f = \begin{pmatrix}0\\0\\ \vdots \\0\\1\end{pmatrix}.
  \label{eq:frobenius_form}
\end{equation}

Преобразование подобия $\mat{Q} = \mat{S}_y\mat{T}$, где
\begin{equation}
  \mat{T} = \begin{pmatrix}
    a_{n-1} & a_{n-2} & \cdots & a_1 & 1\\
    a_{n-2} & a_{n-3} & \cdots & 1 & 0\\
    \vdots & & & & \vdots \\
    1 & 0 & \cdots & & 0
  \end{pmatrix},
  \label{eq:T_matrix}
\end{equation}
$a_i$ --- коэффициенты характеристического полинома исходной матрицы
$\chi_n(\lambda) = \det(\lambda \mat{I} - \mat{A}) = \lambda^n + a_1\lambda^{n-1}+\ldots+a_n$.

Параметры регулятора в исходном базисе (формула 3.4 учебника):
\begin{equation}
  \boxed{\mat{K} = \big(\mat{Q}^{\T}\big)^{-1}(\vect{a} - \vect{a}_{\text{уст}}),}
  \label{eq:modal_gain}
\end{equation}
где $\vect{a}_{\text{уст}} = (a_n^{\text{уст}}, a_{n-1}^{\text{уст}}, \ldots, a_1^{\text{уст}})\T$ ---
коэффициенты желаемого характеристического полинома
$\prod_{i=1}^n(\lambda - \lambda_i^{\text{уст}})$.

\subsection{Многомерный вход}

При $r > 1$ (как у нашего робота: $u_v$ и $u_\omega$) канонизация по
\eqref{eq:frobenius_form} напрямую неприменима. На практике используют
функцию \texttt{place(A,B,p)} пакета Control System Toolbox, реализующую
алгоритм Tits--Yang (численно устойчивый). Для нашей задачи в MATLAB:
\begin{filebox}[title={\filepath{matlab/design\_modal.m}}]
\begin{lstlisting}[style=matlabstyle]
poles_des = [-2; -3; -4; -8; -10];   % желаемые полюса
K_modal   = place(Ad, Bd, exp(poles_des * Ts));   % дискретный
\end{lstlisting}
\end{filebox}

\section{Оптимальный регулятор (ЛКР)}
\label{sec:lqr}

Подход модального синтеза требует ручного выбора полюсов, что
не масштабируется при многомерном состоянии. \textbf{Линейно-квадратичный
регулятор} (ЛКР) автоматически минимизирует интегральный критерий качества
(п.~3.3 учебника):

\begin{formulabox}[title={Критерий качества}]
\begin{equation}
  J = \int_0^\infty \big(\,\vect{x}\T\,\mat{Q}\,\vect{x} + \vect{u}\T\,\mat{R}\,\vect{u}\,\big)\,dt,
  \qquad \mat{Q} = \mat{Q}\T \succeq 0, \quad \mat{R} = \mat{R}\T \succ 0.
  \label{eq:lqr_cost}
\end{equation}
\end{formulabox}

Здесь $\mat{Q}$ штрафует отклонение состояния (точность),
$\mat{R}$ --- расход управления (энергию). Подбор соотношения $\mat{Q}/\mat{R}$
определяет компромисс «точность $\leftrightarrow$ затраты».

\subsection{Уравнение Риккати}

Оптимальное управление (формула 3.10 учебника) имеет вид
\begin{equation}
  \vect{u}^*(t) = -\mat{R}^{-1}\mat{B}\T\mat{P}\,\vect{x}(t) = -\mat{K}\,\vect{x}(t),
  \label{eq:lqr_optimal}
\end{equation}
где $\mat{P}$ --- симметричная положительно определённая матрица,
удовлетворяющая алгебраическому уравнению Риккати (CARE):

\begin{formulabox}[title={Непрерывное уравнение Риккати}]
\begin{equation}
  \mat{A}\T\mat{P} + \mat{P}\mat{A} - \mat{P}\mat{B}\mat{R}^{-1}\mat{B}\T\mat{P} = -\mat{Q}.
  \label{eq:care}
\end{equation}
\end{formulabox}

\textbf{Дискретный аналог} (DARE):
\begin{equation}
  \mat{A}\T\mat{P}\mat{A} - \mat{P} - \mat{A}\T\mat{P}\mat{B}\bar{\mat{R}}^{-1}\mat{B}\T\mat{P}\mat{A} = -\mat{Q},
  \quad \bar{\mat{R}} = \mat{R} + \mat{B}\T\mat{P}\mat{B},
  \label{eq:dare}
\end{equation}
матрица регулятора $\mat{K}_d = \bar{\mat{R}}^{-1}\mat{B}\T\mat{P}\mat{A}_d$.

\subsection{Условия существования решения}

По п.~3.3 учебника, единственное положительно определённое $\mat{P}$
существует при выполнении трёх условий:
\begin{enumerate}
  \item Объект $(\mat{A},\mat{B})$ управляем
        ($\operatorname{rank}\mat{S}_y = n$);
  \item $\mat{R} \succ 0$;
  \item $\mat{Q} \succ 0$, либо $\mat{Q} = \mat{H}\T\mat{H}$ при наблюдаемой паре $(\mat{A}\T, \mat{H}\T)$.
\end{enumerate}

Для робота все три условия выполнены при $v_0 > 0$.

\subsection{Выбор весовых матриц}

Согласно рекомендации [13] из учебника (п.~3.3 после формулы 3.12),
естественный выбор:
\begin{equation}
  \mat{R} = \mat{I}_r,
  \qquad
  \tilde{\mat{Q}} = \operatorname{diag}\!\left(\tfrac{1}{(y^{\text{ст}})^2},\,\tfrac{1}{(t_{\max}/3)^2}\right),
  \label{eq:weight_choice}
\end{equation}
где $y^{\text{ст}}$ --- допустимая статическая ошибка, $t_{\max}$ ---
требуемая длительность переходного процесса.

\textbf{Для робота SAMURAI} разумные параметры:
\begin{equation}
  \mat{Q} = \operatorname{diag}(10,\,10,\,5,\,1,\,1),
  \qquad
  \mat{R} = \operatorname{diag}(1,\,1).
  \label{eq:Q_R_robot}
\end{equation}
Большие веса на $p_x, p_y, \theta$ означают приоритет точности позиции
и курса; малые веса на $v, \omega$ --- допускаем заметные отклонения
скоростей. $\mat{R} = \mat{I}$ --- штрафуем затраты управления равномерно.

\begin{notebox}
\textbf{«Сделать поведение робота ошибочным»} в терминах ЛКР означает:
\begin{itemize}
  \item Заменить $\mat{Q}$ на отрицательно определённую --- алгоритм
        не сходится или даст нестабильный регулятор;
  \item Поменять знак $\mat{R}$ или сделать $\mat{R}$ очень малой ---
        получим избыточно агрессивное управление с насыщением и
        автоколебаниями;
  \item Намеренно использовать желаемые полюса в правой полуплоскости
        $\lambda^{\text{уст}}_i > 0$ для модального регулятора ---
        замкнутая система будет неустойчива.
\end{itemize}
Эти эксперименты безопасны в симуляторе, но опасны на реальном железе ---
проводить только с поднятыми колёсами.
\end{notebox}

\section{Регулирование с ненулевой уставкой}
\label{sec:setpoint}

ЛКР стабилизирует систему в нулевой точке. Для отслеживания
заданной траектории $\vect{x}_0 \neq \vect{0}$ (например, цели Nav2)
требуется дополнительное управляющее воздействие (п.~3.2 учебника,
формула 3.7):
\begin{equation}
  \vect{u}(t) = -\mat{K}\,\vect{x}(t) + \vect{u}_{\text{уст}}(t),
  \quad
  \vect{u}_{\text{уст}} = -\big(\mat{C}\mat{A}_z^{-1}\mat{B}\big)^{-1}\,\vect{y}_{\text{уст}}(t),
  \label{eq:setpoint_feedforward}
\end{equation}
где $\mat{A}_z = \mat{A} - \mat{B}\mat{K}$ --- матрица замкнутой системы.
Это \textbf{прямая связь по уставке} (feedforward), компенсирующая статическую
ошибку при ступенчатом изменении задания.

Для робота уставка $\vect{y}_{\text{уст}} = (p_x^*, p_y^*, \theta^*, v^*, 0)\T$
формируется из заданной точки и желаемой скорости подхода.

\section{Наблюдатель состояния (Луенбергер)}
\label{sec:observer}

Регулятор \eqref{eq:state_feedback} использует полный вектор $\vect{x}(t)$,
но реальные датчики дают только $\vect{y} = \mat{C}\vect{x}$. Согласно
п.~3.4.1 учебника, состояние восстанавливается полноразмерным наблюдателем:

\begin{formulabox}[title={Уравнение наблюдателя (Luenberger)}]
\begin{equation}
  \dot{\hat{\vect{x}}}(t) = \mat{A}\,\hat{\vect{x}} + \mat{B}\,\vect{u}
       + \mat{L}\big(\vect{y} - \mat{C}\hat{\vect{x}}\big).
  \label{eq:observer}
\end{equation}
\end{formulabox}

Ошибка восстановления $\vect{e} = \vect{x} - \hat{\vect{x}}$ удовлетворяет:
\begin{equation}
  \dot{\vect{e}} = (\mat{A} - \mat{L}\mat{C})\,\vect{e}.
  \label{eq:observer_error}
\end{equation}

Матрица $\mat{L}$ выбирается так, чтобы $(\mat{A}-\mat{L}\mat{C})$
имела заданные собственные значения, обеспечивающие сходимость
$\vect{e}(t) \to 0$ существенно быстрее динамики регулятора.
Стандартная рекомендация (п.~3.4.2 учебника): полюса наблюдателя в $5\ldots10$ раз
быстрее полюсов регулятора:
\begin{equation}
  \tilde{\lambda}^{\text{набл}}_i = 5 \cdot \lambda^{\text{уст}}_i.
  \label{eq:observer_eigvals}
\end{equation}

\subsection{Принцип разделения}

При совместном использовании регулятора и наблюдателя замкнутая
система имеет блочную матрицу (формула 3.23 учебника):
\begin{equation}
  \mat{A}_c = \begin{pmatrix}
    \mat{A}+\mat{B}\mat{K} & -\mat{B}\mat{K}\\
    \mat{0} & \mat{A}-\mat{L}\mat{C}
  \end{pmatrix}.
  \label{eq:separation}
\end{equation}
Спектр $\sigma(\mat{A}_c) = \sigma(\mat{A}+\mat{B}\mat{K}) \cup \sigma(\mat{A}-\mat{L}\mat{C})$ ---
полюса регулятора и наблюдателя независимы (\textbf{принцип разделения}).
Это позволяет проектировать $\mat{K}$ и $\mat{L}$ \textbf{независимо}.

\subsection{Условие наблюдаемости}

Существование $\mat{L}$, обеспечивающего любые желаемые собственные
значения, требует полной наблюдаемости (п.~3.4.1):
\begin{equation}
  \operatorname{rank}\mat{S}_n = n,
  \quad
  \mat{S}_n = \begin{pmatrix} \mat{C}\T & \mat{A}\T\mat{C}\T & \cdots & (\mat{A}\T)^{n-1}\mat{C}\T \end{pmatrix}.
  \label{eq:observability}
\end{equation}
При $\mat{C} = \mat{I}$ (как в нашей модели --- все 5 координат
наблюдаемы) это условие выполнено тривиально.

\section{Численный пример для робота}
\label{sec:numerical_example}

Применим теорию к модели \eqref{eq:AB_numerical} при $v_0 = 0.2$ м/с,
$T_s = 0.05$ с. Дискретизация ZOH даёт:
\begin{equation*}
  \mat{A}_d \approx \begin{pmatrix}
    1 & 0 & 0 & 0.0488 & 0\\
    0 & 1 & 0.01 & 0 & 0.000244\\
    0 & 0 & 1 & 0 & 0.0488\\
    0 & 0 & 0 & 0.717 & 0\\
    0 & 0 & 0 & 0 & 0.607
  \end{pmatrix}.
\end{equation*}

Решение DARE с $\mat{Q} = \operatorname{diag}(10,10,5,1,1)$, $\mat{R} = \mat{I}_2$
(вычисление в \filepath{matlab/design\_lqr.m}, см. главу~\ref{ch:implementation})
даёт регулятор:
\begin{equation}
  \mat{K}_{\text{LQR}} \approx \begin{pmatrix}
    2.84 & 0 & 0 & 1.21 & 0\\
    0 & 0.95 & 1.62 & 0 & 0.74
  \end{pmatrix}.
  \label{eq:K_lqr_robot}
\end{equation}
Структура $\mat{K}$ отражает связи: $u_v$ зависит от $p_x$ и $v$
(продольный канал), $u_\omega$ --- от $p_y, \theta, \omega$
(боковой канал, через рыскание).

Собственные значения замкнутой непрерывной системы $\mat{A} - \mat{B}\mat{K}$
лежат в левой полуплоскости (асимптотическая устойчивость), модули собственных
значений $|\lambda_i|$ соответствуют постоянным времени $0{,}3\ldots 1$ с,
что обеспечивает время регулирования $\sim 1{,}5$ с.

\begin{notebox}
ЛКР даёт автоматический компромисс «точность/затраты» через выбор
$\mat{Q},\mat{R}$. Для расширения на \textbf{ограниченные} управления
(физические лимиты: $|u_v| \leq 0{,}3$ м/с, $|u_\omega| \leq 2$ рад/с)
переходим к MPC --- следующая глава.
\end{notebox}

% ============================================================
%  Глава 3 — Model Predictive Control (MPC)
% ============================================================
\chapter{Прогнозирующее управление (MPC)}
\label{ch:mpc}

ЛКР, рассмотренный в главе~\ref{ch:synthesis}, даёт оптимальный регулятор
\textbf{без учёта ограничений} на управление. Реальный робот SAMURAI имеет
жёсткие физические лимиты:
\begin{equation}
  |u_v| \leq u_v^{\max} = 0{,}3 \text{ м/с}, \quad
  |u_\omega| \leq u_\omega^{\max} = 2{,}0 \text{ рад/с}.
  \label{eq:robot_limits}
\end{equation}
Если ЛКР рассчитывает $u_v = 0{,}5$ м/с, актуатор насытит на $0{,}3$ ---
возникает анти-винд-ап эффект, и оптимальность теряется.

\textbf{Model Predictive Control} (MPC) решает эту задачу: на каждом
шаге решается \emph{оптимизационная задача} с явным учётом ограничений
на горизонте $N$ шагов вперёд.

\section{Идея метода}
\label{sec:mpc_idea}

В момент времени $t = kT_s$:
\begin{enumerate}
  \item Измеряется текущее состояние $\vect{x}_k$ (или его оценка
        $\hat{\vect{x}}_k$ от наблюдателя).
  \item Решается задача оптимизации: найти последовательность
        $\vect{u}_k, \vect{u}_{k+1}, \ldots, \vect{u}_{k+N-1}$,
        минимизирующую функционал на горизонте.
  \item Применяется только \emph{первый} элемент $\vect{u}_k$.
  \item На следующем шаге $k+1$ горизонт сдвигается ---
        \textbf{принцип сдвигаемого горизонта} (receding horizon).
\end{enumerate}

\begin{notebox}
MPC можно рассматривать как ЛКР с конечным горизонтом плюс ограничения.
При $N \to \infty$ и без ограничений MPC сводится к ЛКР с тем же
матрицей $\mat{K}$.
\end{notebox}

\section{Математическая постановка}
\label{sec:mpc_formulation}

Для дискретной модели \eqref{eq:discrete_model}:

\begin{formulabox}[title={Задача MPC на горизонте $N$}]
\begin{align}
  \min_{\{\vect{u}_i\}_{i=k}^{k+N-1}}
  &\quad J = \sum_{i=k}^{k+N-1}\!\!\big(\vect{x}_i\T\mat{Q}\vect{x}_i + \vect{u}_i\T\mat{R}\vect{u}_i\big)
            + \vect{x}_{k+N}\T\mat{P}_f\,\vect{x}_{k+N}
  \label{eq:mpc_cost}\\[4pt]
  \text{при условиях:}
  &\quad \vect{x}_{i+1} = \mat{A}_d\vect{x}_i + \mat{B}_d\vect{u}_i,
  \quad i = k, \ldots, k+N-1,
  \label{eq:mpc_dynamics}\\
  &\quad \vect{u}_{\min} \leq \vect{u}_i \leq \vect{u}_{\max},
  \label{eq:mpc_input_constr}\\
  &\quad \vect{x}_{\min} \leq \vect{x}_i \leq \vect{x}_{\max}.
  \label{eq:mpc_state_constr}
\end{align}
\end{formulabox}

Здесь $\mat{P}_f$ --- терминальный штраф (terminal cost), обычно
выбирается как решение уравнения Риккати $\mat{P}_\infty$ ---
это гарантирует устойчивость замкнутой системы при $N$ достаточно большом.

\section{Сведение к QP}
\label{sec:mpc_qp}

Задача \eqref{eq:mpc_cost}--\eqref{eq:mpc_state_constr} ---
\textbf{квадратичная программа} (QP) по переменной
$\vect{U} = (\vect{u}_k\T, \ldots, \vect{u}_{k+N-1}\T)\T \in \mathbb{R}^{rN}$.

\subsection{Развёртка динамики}

Подставив \eqref{eq:mpc_dynamics} итеративно, выразим все будущие состояния
через текущее $\vect{x}_k$ и $\vect{U}$:
\begin{equation}
  \begin{pmatrix}\vect{x}_{k+1}\\ \vect{x}_{k+2}\\ \vdots\\ \vect{x}_{k+N}\end{pmatrix}
  =
  \underbrace{\begin{pmatrix}\mat{A}_d\\ \mat{A}_d^2\\ \vdots\\ \mat{A}_d^N\end{pmatrix}}_{\bm{\Phi}}\,\vect{x}_k
  +
  \underbrace{\begin{pmatrix}
    \mat{B}_d & 0 & \cdots & 0\\
    \mat{A}_d\mat{B}_d & \mat{B}_d & \cdots & 0\\
    \vdots & \vdots & \ddots & \vdots\\
    \mat{A}_d^{N-1}\mat{B}_d & \mat{A}_d^{N-2}\mat{B}_d & \cdots & \mat{B}_d
  \end{pmatrix}}_{\bm{\Gamma}}
  \,\vect{U}.
  \label{eq:lifting}
\end{equation}

\subsection{QP в стандартной форме}

Подставим \eqref{eq:lifting} в \eqref{eq:mpc_cost}, получим:

\begin{formulabox}[title={Стандартная QP-задача MPC}]
\begin{equation}
  \min_{\vect{U}} \;\tfrac{1}{2}\vect{U}\T \mat{H}\,\vect{U} + \vect{f}\T(\vect{x}_k)\,\vect{U}
  \quad \text{при} \quad
  \mat{G}\vect{U} \leq \vect{g}(\vect{x}_k),
  \label{eq:qp_standard}
\end{equation}
\end{formulabox}

где
\begin{equation*}
  \mat{H} = 2\big(\bm{\Gamma}\T\bar{\mat{Q}}\bm{\Gamma} + \bar{\mat{R}}\big),
  \qquad
  \vect{f}(\vect{x}_k) = 2\bm{\Gamma}\T\bar{\mat{Q}}\bm{\Phi}\,\vect{x}_k,
\end{equation*}
$\bar{\mat{Q}} = \operatorname{blkdiag}(\mat{Q}, \ldots, \mat{Q}, \mat{P}_f)$,
$\bar{\mat{R}} = \operatorname{blkdiag}(\mat{R}, \ldots, \mat{R})$.

Матрица $\mat{H}$ постоянна (не зависит от $\vect{x}_k$) и может быть
вычислена offline. Вектор $\vect{f}$ и правая часть $\vect{g}$
ограничений пересчитываются на каждом шаге.

\section{Алгоритм решения}
\label{sec:mpc_algorithm}

Для робота SAMURAI используется простой и надёжный подход:

\subsection{Без ограничений: явное решение}

Если ограничения \eqref{eq:mpc_input_constr} неактивны
(управление лежит внутри допустимой области), QP \eqref{eq:qp_standard}
имеет аналитическое решение:
\begin{equation}
  \vect{U}^* = -\mat{H}^{-1}\vect{f}(\vect{x}_k) = -\mat{K}_{\text{MPC}}\,\vect{x}_k,
  \label{eq:mpc_unconstrained}
\end{equation}
где $\mat{K}_{\text{MPC}} \in \mathbb{R}^{rN \times n}$ ---
\emph{статический} коэффициент, вычисляемый один раз. Применяется первая
строка размера $r \times n$ (первое управление $\vect{u}_k$).

\textbf{Это ровно та же $\mat{K}$, что и у ЛКР, при $N\to\infty$ и
$\mat{P}_f = \mat{P}_\infty$.} Для $N = 10\ldots 30$ результат
близок к ЛКР, но позволяет ввести ограничения.

\subsection{С ограничениями: QP-солвер}

При активных ограничениях задача \eqref{eq:qp_standard} решается
итеративно. В нашей реализации (см. главу~\ref{ch:implementation})
используется простой \textbf{проекционный метод}: после явного решения
\eqref{eq:mpc_unconstrained} результат проецируется на бокс
$[\vect{u}_{\min}, \vect{u}_{\max}]$:
\begin{equation}
  \vect{u}_k^{\text{clip}} = \max\!\big(\vect{u}_{\min},\,\min(\vect{u}_{\max}, \vect{u}_k^*)\big).
  \label{eq:projection}
\end{equation}

\begin{notebox}
Это \textbf{не оптимальное} решение QP с ограничениями (для оптимальности
нужен active-set или interior-point), но даёт корректное и устойчивое
поведение для робота. Полноценный QP-солвер (\texttt{cvxpy}, \texttt{qpOASES},
\texttt{scipy.optimize.minimize}) можно подключить опцией
\texttt{control.mpc.solver: qp} в конфиге.
\end{notebox}

\subsection{Терминальный штраф $\mat{P}_f$}

Согласно теореме об устойчивости MPC (Mayne et al., 2000), терминальный
штраф $\mat{P}_f$, выбранный как решение уравнения Риккати замкнутой
системы $\mat{A}_d - \mat{B}_d\mat{K}_\infty$, гарантирует устойчивость
MPC при любом $N \geq 1$. Поэтому в \filepath{matlab/design\_mpc.m}
сначала решается DARE для получения $\mat{P}_f = \mat{P}_\infty$,
затем формируется QP.

\section{Параметры MPC для SAMURAI}
\label{sec:mpc_params}

\begin{center}
\begin{tabular}{lll}
\toprule
\textbf{Параметр} & \textbf{Значение} & \textbf{Комментарий}\\
\midrule
$T_s$ & $0{,}05$ с & период цикла \texttt{motor\_node} (20 Гц)\\
$N$ & $10$ & горизонт $0{,}5$ с (соизмерим с $\tau_v$)\\
$\mat{Q}$ & $\operatorname{diag}(10,10,5,1,1)$ & точность позиции/курса\\
$\mat{R}$ & $\operatorname{diag}(1,1)$ & умеренный штраф управления\\
$\mat{P}_f$ & решение DARE & гарантия устойчивости\\
$\vect{u}_{\min}$ & $(-0{,}3,\,-2{,}0)$ & физические лимиты робота\\
$\vect{u}_{\max}$ & $(+0{,}3,\,+2{,}0)$ & --- \texttt{simulator.max\_*}\\
\bottomrule
\end{tabular}
\end{center}

При $N = 10$, $n = 5$, $r = 2$ задача QP имеет $20$ переменных,
$40$ ограничений-неравенств. Время решения на Raspberry Pi 4 ---
менее 1 мс (запас порядка для 50 мс цикла).

\section{Подбор поведения и «ошибочные» режимы}
\label{sec:behavior_tuning}

Параметры MPC вынесены в \filepath{config.yaml} и могут быть изменены
без перекомпиляции. Это позволяет:

\paragraph{1. Корректные поведения}
\begin{itemize}
  \item \textbf{Точный медленный режим}: $\mat{Q}\!\uparrow$,
        $\mat{R}\!\uparrow$, $u_{\max}\!\downarrow$ ---
        робот ездит медленно, но точно.
  \item \textbf{Агрессивный режим}: $\mat{Q}_{p_x,p_y}\!\uparrow\uparrow$,
        $\mat{R}\!\downarrow$ --- быстрая отработка, рывки.
  \item \textbf{Длинный горизонт}: $N = 30$ --- видит дальше, плавнее
        огибает препятствия (но дороже по CPU).
\end{itemize}

\paragraph{2. «Ошибочные» поведения для отладки}
\begin{itemize}
  \item \textbf{Отрицательный штраф}: $\mat{Q} = -\mat{I}$ ---
        QP станет невыпуклой, регулятор будет «гнать» в нестабильность.
  \item \textbf{Завышенные лимиты}: $u_{\max} = 100$ --- ЛКР перестанет
        насыщаться, проявится колебание из-за неточностей модели.
  \item \textbf{Полюса в правой полуплоскости} (модальный режим):
        $\lambda^{\text{уст}} = +1$ --- замкнутая система неустойчива,
        робот «убегает».
  \item \textbf{Зануление обратной связи по позиции}: $\mat{Q}_{1,1}=\mat{Q}_{2,2}=0$ ---
        робот игнорирует координаты $p_x, p_y$, движется только по
        внутренним скоростям.
\end{itemize}

\begin{warnbox}
Все «ошибочные» режимы безопасны в симуляторе (\texttt{./samurai.sh sim}),
но опасны на реальном железе. Для физических испытаний:
\begin{itemize}
  \item Поднимайте робота на подставке --- колёса не должны касаться пола;
  \item Включите аппаратный аварийный стоп
        (\texttt{config.motor.collision\_guard\_stop\_m});
  \item Не превышайте $u_{\max} = 0{,}5$ м/с.
\end{itemize}
\end{warnbox}

% ============================================================
%  Глава 4 — Реализация регулятора в роботе
% ============================================================
\chapter{Реализация в коде робота}
\label{ch:implementation}

В этой главе описывается \textbf{полный конвейер}:
\textsf{MATLAB (offline)} $\to$ \textsf{YAML/JSON} $\to$ \textsf{Python (online)}
$\to$ \textsf{motor\_node} реального робота. Все коэффициенты выносятся
в \filepath{config.yaml}, что позволяет менять поведение робота
без перекомпиляции.

\section{Архитектура}
\label{sec:architecture}

\begin{center}
\begin{tikzpicture}[
  box/.style={draw, rounded corners, minimum width=3cm, minimum height=0.9cm,
              font=\small, align=center},
  arrow/.style={-{Stealth[length=2.5mm]}, thick},
  node distance=0.6cm and 1.5cm
]
\node[box, fill=blue!8]   (matlab)  {MATLAB\\\filepath{matlab/main.m}};
\node[box, fill=green!8, right=of matlab] (yaml) {\filepath{config.yaml}\\\small \texttt{control:} раздел};
\node[box, fill=orange!8, right=of yaml] (py)   {Python\\\filepath{pi\_nodes/control/}};
\node[box, fill=red!8, below=of py] (motor) {\filepath{motor\_node.py}\\цикл 20 Гц};

\draw[arrow] (matlab) -- node[above, font=\scriptsize] {export} (yaml);
\draw[arrow] (yaml)   -- node[above, font=\scriptsize] {load} (py);
\draw[arrow] (py)     -- node[right, font=\scriptsize] {cmd\_vel} (motor);
\end{tikzpicture}
\end{center}

\textbf{Шаги работы:}
\begin{enumerate}
  \item Инженер запускает \texttt{main.m} в MATLAB --- получает $\mat{A}, \mat{B}, \mat{K}, \mat{L}, \mat{P}_f$;
  \item Скрипт автоматически экспортирует матрицы в \filepath{config.yaml}
        раздел \texttt{control};
  \item Робот при старте читает \texttt{control} и инициализирует регулятор;
  \item На каждом цикле \texttt{motor\_node} вызывает $\vect{u}_k = \texttt{controller.step}(\vect{x}_k, \vect{x}_k^{\text{ref}})$;
  \item Если \texttt{control.mode} = \texttt{off}, действует старая логика
        прямой трансляции \texttt{cmd\_vel}.
\end{enumerate}

\section{Конфигурация}
\label{sec:config}

Все коэффициенты вынесены в \filepath{config.yaml}, раздел \texttt{control}:

\begin{filebox}[title={\filepath{config.yaml} --- \texttt{control:}}]
\begin{lstlisting}[basicstyle=\footnotesize\ttfamily]
control:
  mode: "off"           # off | lqr | mpc | modal — режим регулятора

  plant:                # параметры объекта (физика)
    v0:        0.20     # м/с — опорная скорость линеаризации
    tau_v:     0.15     # с — постоянная мотора (linear)
    tau_w:     0.10     # с — постоянная мотора (angular)
    Ts:        0.05     # с — период дискретизации (= 1 / 20 Hz)

  weights:              # критерий качества
    Q_diag:  [10, 10, 5, 1, 1]   # x,y,theta,v,omega
    R_diag:  [1, 1]              # u_v, u_omega

  limits:               # физические ограничения управления
    u_min:   [-0.30, -2.0]
    u_max:   [+0.30, +2.0]

  mpc:                  # параметры MPC
    horizon_N:  10      # горизонт прогноза (шагов)
    solver:     "clip"  # clip | qp (qp требует scipy/cvxpy)

  modal:                # параметры модального регулятора
    poles_continuous: [-2, -3, -4, -8, -10]   # желаемые
    observer_speedup: 5.0   # во сколько раз быстрее регулятора

  # Матрицы рассчитаны MATLAB-скриптом, обновляются автоматически.
  # Для ручной правки: установите override: true и поменяйте значения.
  matrices:
    override: false
    A: null   # 5x5 — заполнится из MATLAB
    B: null   # 5x2
    K: null   # 2x5 — главный регулятор
    L: null   # 5x5 — наблюдатель (если включён)
    Pf: null  # 5x5 — терминальный штраф MPC
\end{lstlisting}
\end{filebox}

\section{Python-модуль управления}
\label{sec:python_impl}

В пакете \filepath{pi\_nodes/control/} добавлены три новых модуля:

\begin{center}
\begin{tabular}{ll}
\toprule
\textbf{Файл} & \textbf{Назначение}\\
\midrule
\filepath{state\_space\_model.py} & Дискретная модель $\vect{x}_{k+1}=\mat{A}_d\vect{x}_k + \mat{B}_d\vect{u}_k$\\
\filepath{lqr\_controller.py}     & ЛКР: $\vect{u}_k = -\mat{K}\,\vect{x}_k$\\
\filepath{mpc\_controller.py}     & MPC: явное решение + проекция на бокс\\
\bottomrule
\end{tabular}
\end{center}

\subsection{Загрузка матриц}

\begin{filebox}[title={\filepath{pi\_nodes/control/state\_space\_model.py}}]
\begin{lstlisting}[style=pythonstyle]
import numpy as np
from config_loader import cfg

class StateSpaceModel:
    """Discrete state-space model x[k+1] = Ad*x[k] + Bd*u[k]."""
    def __init__(self, Ad=None, Bd=None):
        if Ad is None:
            Ad = np.array(cfg('control.matrices.A'), dtype=float)
        if Bd is None:
            Bd = np.array(cfg('control.matrices.B'), dtype=float)
        self.Ad, self.Bd = Ad, Bd
        self.n, self.r = self.Bd.shape

    def step(self, x, u):
        return self.Ad @ x + self.Bd @ u
\end{lstlisting}
\end{filebox}

\subsection{Регулятор LQR}

\begin{filebox}[title={\filepath{pi\_nodes/control/lqr\_controller.py}}]
\begin{lstlisting}[style=pythonstyle]
import numpy as np
from config_loader import cfg

class LQRController:
    def __init__(self, K=None, u_min=None, u_max=None):
        self.K     = np.array(K     if K     is not None
                              else cfg('control.matrices.K'))
        self.u_min = np.array(u_min if u_min is not None
                              else cfg('control.limits.u_min'))
        self.u_max = np.array(u_max if u_max is not None
                              else cfg('control.limits.u_max'))

    def step(self, x, x_ref=None):
        dx = x - x_ref if x_ref is not None else x
        u  = -self.K @ dx
        return np.clip(u, self.u_min, self.u_max)
\end{lstlisting}
\end{filebox}

\subsection{MPC: явное решение + проекция}

\begin{filebox}[title={\filepath{pi\_nodes/control/mpc\_controller.py} (упрощённо)}]
\begin{lstlisting}[style=pythonstyle]
class MPCController:
    """MPC c явным решением + проекцией на бокс ограничений."""

    def __init__(self, Ad, Bd, Q, R, Pf, N, u_min, u_max):
        self.N = N
        self._build_qp_matrices(Ad, Bd, Q, R, Pf)
        self.u_min, self.u_max = u_min, u_max
        self.r = Bd.shape[1]

    def _build_qp_matrices(self, Ad, Bd, Q, R, Pf):
        n, r = Bd.shape
        N = self.N
        # Phi (n*N x n) и Gamma (n*N x r*N) -- развёртка динамики
        Phi   = np.vstack([np.linalg.matrix_power(Ad, i+1) for i in range(N)])
        Gamma = np.zeros((n*N, r*N))
        for i in range(N):
            for j in range(i+1):
                block = np.linalg.matrix_power(Ad, i-j) @ Bd
                Gamma[i*n:(i+1)*n, j*r:(j+1)*r] = block

        Q_bar = np.kron(np.eye(N-1), Q)
        Q_bar = block_diag_pad(Q_bar, Pf)            # терминальный штраф
        R_bar = np.kron(np.eye(N), R)

        H_full = 2*(Gamma.T @ Q_bar @ Gamma + R_bar)
        f_mat  = 2*Gamma.T @ Q_bar @ Phi

        # Только первое управление: K_first = (H^-1 f)[:r, :]
        K_first = -np.linalg.solve(H_full, f_mat)[:r, :]
        self.K_first = K_first

    def step(self, x, x_ref=None):
        dx = x - x_ref if x_ref is not None else x
        u  = self.K_first @ dx
        return np.clip(u, self.u_min, self.u_max)
\end{lstlisting}
\end{filebox}

\section{Тестирование}
\label{sec:testing}

Для каждого модуля написан unit-тест в \filepath{tests/}.

\subsection{test\_state\_space\_model.py}

Проверяет:
\begin{itemize}
  \item $\vect{x}_{k+1} = \mat{A}_d\vect{x}_k + \mat{B}_d\vect{u}_k$ ---
        корректность дискретного шага;
  \item Сходимость к точке равновесия при $\vect{u} = \vect{0}$
        (для устойчивого $\mat{A}_d$).
\end{itemize}

\subsection{test\_lqr\_controller.py}

Проверяет:
\begin{itemize}
  \item Замкнутая система $\mat{A}_d - \mat{B}_d\mat{K}$ устойчива
        ($|\lambda_i| < 1$);
  \item Управление лежит в пределах $[\vect{u}_{\min}, \vect{u}_{\max}]$;
  \item При $\vect{x} = \vect{0}$ управление равно нулю.
\end{itemize}

\subsection{test\_mpc\_controller.py}

Проверяет:
\begin{itemize}
  \item Без ограничений MPC совпадает с ЛКР для горизонта $N \geq 20$;
  \item При активных ограничениях управление лежит на границе бокса;
  \item Замкнутая симуляция сходится к нулю.
\end{itemize}

\section{Примеры использования}
\label{sec:examples}

\subsection{Запуск симуляции с MPC}

\begin{filebox}[title={Изменение режима регулятора}]
\begin{lstlisting}[basicstyle=\small\ttfamily]
# 1. В config.yaml:
control:
  mode: mpc

# 2. Запустить симулятор:
./samurai.sh sim

# 3. Открыть dashboard, отправить goal_pose
\end{lstlisting}
\end{filebox}

\subsection{Расчёт нового регулятора в MATLAB}

\begin{filebox}[title={\filepath{matlab/main.m} --- сценарий}]
\begin{lstlisting}[style=matlabstyle]
% Загрузить параметры объекта
params = samurai_params();
% Линеаризовать
[A, B, C, D] = linearize_samurai(params);
% Дискретизировать
[Ad, Bd]     = discretize_samurai(A, B, params.Ts);
% Синтез ЛКР
[K_lqr, P]   = design_lqr(Ad, Bd, params.Q, params.R);
% Синтез MPC
mpc          = design_mpc(Ad, Bd, params.Q, params.R, P, params.N);
% Экспорт в config.yaml
export_to_yaml(Ad, Bd, K_lqr, P, mpc, '../config.yaml');
% Симуляция
simulate_closed_loop(Ad, Bd, K_lqr);
\end{lstlisting}
\end{filebox}

\subsection{«Ошибочный» режим для демонстрации}

Чтобы показать студентам \textbf{неустойчивую систему}, достаточно
поставить в \texttt{config.yaml}:
\begin{filebox}[title={Демонстрация неустойчивости}]
\begin{lstlisting}[basicstyle=\small\ttfamily]
control:
  mode: modal
  modal:
    poles_continuous: [+1.0, +0.5, -1, -1, -1]   # 2 положительных!
\end{lstlisting}
\end{filebox}

При запуске симуляции робот будет «убегать» по экспоненте,
так как замкнутая система имеет полюса в правой полуплоскости.

\section{Заключение}
\label{sec:conclusion}

\textbf{Резюме созданной системы:}

\begin{enumerate}
  \item Построена 5-мерная нелинейная ОДУ-модель робота
        (глава~\ref{ch:model});
  \item Модель линеаризована вокруг режима прямолинейного движения
        и дискретизирована с $T_s = 0{,}05$ с;
  \item Синтезированы 3 регулятора: модальный (pole placement),
        ЛКР (Risk-метод), MPC с горизонтом $N = 10$ и проекцией;
  \item Полноразмерный наблюдатель восстанавливает все 5 координат
        состояния по доступным измерениям;
  \item Расчёт выполняется offline в MATLAB
        (\filepath{matlab/}, 6 файлов), результат экспортируется в
        \filepath{config.yaml};
  \item Python-реализация (\filepath{pi\_nodes/control/}, 3 файла)
        читает конфиг при старте и применяет регулятор в цикле
        \texttt{motor\_node} (20 Гц);
  \item Покрыто unit-тестами (\filepath{tests/test\_*\_controller.py}, 3 файла);
  \item Все параметры $\mat{Q}, \mat{R}, N, \vect{u}_{\min,\max}$ ---
        редактируемы из конфига без перекомпиляции, что позволяет как
        улучшать поведение робота, так и намеренно вводить ошибки
        для отладки и обучения.
\end{enumerate}

\textbf{Ссылки на файлы проекта:}

\begin{center}
\begin{tabular}{ll}
\toprule
\textbf{Файл} & \textbf{Содержимое}\\
\midrule
\filepath{matlab/main.m}                                & оркестратор расчёта\\
\filepath{matlab/samurai\_params.m}                     & параметры объекта\\
\filepath{matlab/linearize\_samurai.m}                  & матрицы $\mat{A}, \mat{B}$\\
\filepath{matlab/discretize\_samurai.m}                 & ZOH-дискретизация\\
\filepath{matlab/design\_lqr.m}                         & решение DARE\\
\filepath{matlab/design\_mpc.m}                         & MPC-матрицы\\
\filepath{matlab/simulate\_closed\_loop.m}              & симуляция\\
\filepath{matlab/export\_to\_yaml.m}                    & запись конфига\\
\filepath{pi\_nodes/control/state\_space\_model.py}     & дискретная модель\\
\filepath{pi\_nodes/control/lqr\_controller.py}         & ЛКР-регулятор\\
\filepath{pi\_nodes/control/mpc\_controller.py}         & MPC-регулятор\\
\filepath{config.yaml}                                  & раздел \texttt{control:}\\
\filepath{tests/test\_state\_space\_model.py}           & тест модели\\
\filepath{tests/test\_lqr\_controller.py}               & тест ЛКР\\
\filepath{tests/test\_mpc\_controller.py}               & тест MPC\\
\bottomrule
\end{tabular}
\end{center}


\end{document}
